\section{Recommended Reading}

\begin{readingbox}
\textbf{Foundation}
\begin{itemize}[nosep]
  \item John Holland, \emph{Emergence}.
  \item \emph{Introduction to the Theory of Complex Systems}.
  \item \emph{Complexity Science}, for a broader systems vocabulary around
        self-organization and multiscale order.
\end{itemize}
\textbf{Biological and Review Layer}
\begin{itemize}[nosep]
  \item \emph{Metastability demystified: the foundational past, the pragmatic
        present and the promising future} (Nature Reviews Neuroscience, 2024,
        medium).
  \item \emph{Associative conditioning in gene regulatory network models
        increases integrative causal emergence} (Communications Biology, 2025,
        medium).
  \item \emph{Cortical connectivity, local dynamics and stability correlates of
        global conscious states} (Communications Biology, 2025, strong as
        biological background, medium for this chapter's transfer claims).
\end{itemize}
\textbf{Mathematical bridge}
\begin{itemize}[nosep]
  \item Eugene Izhikevich, \emph{Dynamical Systems in Neuroscience}.
  \item Review literature on attractor and integrator networks in the brain
        (Nature Reviews Neuroscience, 2022, medium).
\end{itemize}
\textbf{Reading discipline}
\begin{itemize}[nosep]
  \item Use the biology and review layer to justify organized state formation
        and stability.
  \item Use the systems and mathematics layer to clarify what kind of higher-
        level pattern counts as emergent.
  \item Do not treat frontier consciousness-physics proposals as evidence for
        this chapter's core claims.
\end{itemize}
\end{readingbox}
